% Conclusion

Biological Petri net theory lacks formal semantics for hierarchical regulatory signals consumed during decision-making, preventing expression of regulatory consumption phenomena essential for commitment finality. We introduce Signal Hierarchy Theory providing signal places ($\Psi$) with consumption semantics, enabling hierarchical preemption where higher regulatory layers control lower metabolic layers. We prove layer assignment consistency and preemption correctness, establishing theoretical foundations. The formalism bridges the expressiveness gap between metabolic mass transfer (normal arcs) and regulatory information flow (signal flow arcs) within unified Bio-PN framework.

\vspace{3pt}
\noindent The formalism enables analytical capabilities previously inexpressible in classical Bio-PN: threshold determination (energy concentrations triggering pathway commitment), basin geometry quantification (cell fate reliability via reachability analysis), and decision finality analysis (irreversible commitment through signal consumption). Application to \textit{B. subtilis} sporulation validates predictive capability: SHYPN analysis computes ATP commitment threshold at 2.38 mM, matching experimental measurement (2.21 mM, 7\% error). This quantitative prediction is impossible in classical Bio-PN lacking consumption semantics. The validated framework transforms regulatory networks from qualitative circuit diagrams into predictive mathematical models.

\vspace{3pt}
\noindent Researchers can now formally pose and computationally answer questions requiring signal consumption semantics: What energy threshold triggers irreversible pathway commitment? How reliable is cell fate convergence within attractor basins? When does regulatory decision become irreversible? Signal hierarchy provides formal semantics and validated computational methods addressing these questions with mathematical precision and experimental accuracy previously unavailable in Bio-PN theory.

\vspace{3pt}
\noindent Applications span synthetic biology (engineered decision circuits with guaranteed commitment properties and quantifiable energy requirements), drug discovery (ATP-dependent transporter hierarchies with computed threshold modulation for combination therapy), and precision medicine (quantitative cell fate prediction under therapeutic intervention validated against sporulation data). The unified framework bridges molecular mechanisms with cellular decision-making, supporting rational design of biological systems with formal guarantees and predictive accuracy.

\vspace{3pt}
\noindent Signal Hierarchy Theory extends Bio-PN formalism with consumption semantics, hierarchical organization, theoretical proofs, and validated predictive capability, capturing both metabolic detail and regulatory control in coherent theoretical framework with demonstrated practical utility. Open-source implementation ensures reproducibility for quantitative systems biology applications.
